\section{Engineering Accreditation}\label{engineering-accreditation}

\stanceinfo{\textbf{CFES Congress 2018} [2018-01-07]}{2024-01-02}

\officialstance{The Canadian Federation of Engineering Students believes that the accreditation system for Canadian engineering programs should protect the interests of engineering students by ensuring high, consistent and achievable standards for the quality of engineering education by involving student voices in the process of accreditation visits and the development of accreditation criteria.}

\stanceheading{The Students' Position}
\begin{itemize}
 \item Students are a key stakeholder in engineering education, and should be actively consulted in discussions and decisions related to engineering accreditation.
 \item The accreditation visit process should be reformed to better evaluate and incorporate student perspectives throughout the entire visit.
 \item An effective accreditation system should be primarily focused on learning outcomes, and any measure of student input should be representative of total learning time, and not favor lecture-based learning over other learning methods.
 \item Accreditation Units (AU) do not always favor engineering students, nor affect the quality of engineering education, where a more holistic approach should be taken either through AU or alternatively a different metric for completion of the engineering attributes.
\end{itemize}

\stanceheading{Recommendations to Partners, Stakeholders, and Other Entities}
\begin{itemize}
 \item The CFES calls on the Canadian Engineering Accreditation Board (CEAB) to update its accreditation visit practices to incorporate more valuable and active student feedback at the end of each accreditation year.
 \item The CFES calls on the CEAB to continue building on efforts such as the Graduate Attributes to better reflect learning outcomes and total learning time.
 \item The CFES calls on the CEAB to revise the Accreditation Units, as education and concepts more than just a quantifiable number or checkbox as they may not accurately depict student understanding.
 \item The CFES calls on the CEAB to include CFES student representatives as equal stakeholders in all relevant projects, committees, and task forces related to accreditation [achieved].
 \item The CFES calls on the Engineering Deans Canada (EDC) to collaborate on joint recommendations to the current accreditation system based on the shared interests of faculty and students.
 \item The CFES calls on the transparent communication of EDC with their respective schools to ensure student voices can be heard at the time of the visit to ensure the best engineering education is provided.
 \item The CFES calls on the CEAB to consider incorporating more possibilities for extracurricular learning experiences into the current accreditation system. For the quality of such learning experiences, CFES will work with the CEAB and provisional engineering regulators to create a rigorous framework and set of standards.
 \item The CFES calls on the CEAB to standardize and regulate the quality of internships/coops offered by university engineering programs in Canada.
 \item The CFES calls on EDC to regularly update their course offerings and policies to match that of the changing AU requirements and improvements every academic year.
\end{itemize}

\stanceheading{What the CFES is doing}
\begin{itemize}
 \item The CFES has provided student representation as an observer of the Canadian Engineering Accreditation Board (CEAB) for several years, and in 2015 formalized the responsibility to advocate to the CEAB.
 \item In 2017, the CFES sent representation to a meeting of the EDC for the first time, and has discussed future collaboration on issues related to accreditation.
 \item CEAB meetings grant the CFES speaking rights during open sessions as described by \href{https://engineerscanada.ca/sites/default/files/2022-02/Board-Policy-Manual-Combined-e.pdf}{Engineers Canada Board Policy} Below [1]:
 \begin{itemize}
  \item 6.9.1 F.1 Observers at Meetings (1) The CEAB shall invite the following representatives to its meetings, as observers, each of whom shall be granted the right to be recognized as a speaker in the CEAB's open sessions: a) The president of the Canadian Federation of Engineering Students (CFES), or the CFES president's designate
  \item In addition, \href{https://engineerscanada.ca/sites/default/files/2022-02/Board-Policy-Manual-Combined-e.pdf}{Engineers Canada's Board Policy 7.2} directs the CEAB to maintain a relationship with CFES: c) Engineering students are a key stakeholder of accreditation. In addition to soliciting student feedback during program visits, the Accreditation Board is directed to maintain a relationship with the CFES and invite a representative to observe their meetings, requesting that they bring a report for the CEAB's consideration. All travel costs for this representative are covered by Engineers Canada.
 \end{itemize}
 \item CEAB has updated its sample accreditation visit schedule based on feedback from stakeholders, including the CFES -- creating more time for student engagement during the accreditation visits.
 \item The CFES has included an engineering curriculum section into the 2023 national survey to gauge the general knowledge of CEAB as well as understanding and inclusion of AU.
 \item Consultation with the CEAB was done for the 2023 national survey questions and updated stance documentation.
\end{itemize}

\stanceheading{What the CFES plans to do}
\begin{itemize}
 \item The CFES will push for inclusion in all relevant projects and committees related to the CEAB.
 \item The CFES will continuously create and distribute a national survey in order to update its knowledge on student experiences with accreditation and accreditation-based services. Survey questions will be consulted with CEAB yearly to ensure no crossover with their own survey.
 \item The CFES will continuously update its stance on accreditation as significant changes in policy and regulations arise.
 \item The CFES will hold EDC accountable for untracked changes within their universities of new policy and regulations specific to engineering attributes.
 \item The CFES will gather additional feedback specifically related to accreditation visits to review the student perspective of the quality of the accreditation visits held.
\end{itemize}

\newpage
